% Responsable(s): por asignar
% Objetivo: describir la revisión sistemática: bases consultadas, cadenas de
% búsqueda, criterios de inclusión/exclusión y proceso de selección de artículos.

% Redacción pendiente.
La investigación se desarrolló mediante una revisión bibliográfica de
publicaciones científicas relacionadas con el tiempo de acceso y el rendimiento
de diferentes tecnologías de memoria. La búsqueda se orientó principalmente a
estudios sobre latencia, ancho de banda, tiempo de acceso, jerarquía de memoria
y tecnologías como DDR4, DDR5 y HBM.

Para la selección de las fuentes se consideraron artículos científicos
publicados entre 2022 y 2026. Se seleccionaron al menos doce artículos
relacionados directamente con el tema, procurando que como mínimo un 60\% de
las publicaciones estuvieran escritas en inglés. Las fuentes fueron localizadas
mediante bases de datos académicas disponibles a través de la Biblioteca del
Tecnológico de Costa Rica, Google académico y IEEEXplore.

Una vez seleccionados los artículos, se identificaron los principales datos
relacionados con la tecnología de memoria, latencia, ancho de
banda, tiempo de acceso, y arquitectura utilizada. Seguidamente, estos datos se 
organizaron y compararon con el propósito de determinar las diferencias de rendimiento 
entre las tecnologías analizadas y de esta manera analizar sus características teóricas y 
su comportamiento en sistemas reales.

Finalmente, los resultados de los diferentes estudios se analizaron de forma
colectiva para analizar limitaciones y avances recientes en los
sistemas de memoria. Este proceso permitió utilizar tanto resultados
cuantitativos como observaciones de los autores para desarrollar el estado del
arte, la comparación de tecnologías y las conclusiones preliminares de la
investigación.
